% MAS1004 Assignment 2: long report.
%
% You fill this in with your agent, in your repository. AGENTS.md tells the
% agent how. Red text in the PDF is what still has to be replaced.
%
% - Pictures go in figures/. Copy them from results/, which is not in git.
% - Outputs go in outputs/: copy experiments.csv, LOG.md and the outputs of
%   check_submission.py and curate.py count here. \outputfile{outputs/...}
%   prints a file exactly as it is. Every number in this report has to come
%   from such an output.
%
% To make the PDF: compress this folder into a zip, and on Overleaf choose
% New Project, Upload Project. Set Menu, Compiler to XeLaTeX, and Recompile.

\documentclass[11pt]{article}
\usepackage[ko]{mas1004}

\reporttitle{Assignment 2: Long report}
\studentname{이천희}
\studentid{20251183}

\begin{document}
\makeheader

리더보드 닉네임: NOWHEREMAN. 마지막 릴리스: Track A \texttt{track-a-8} (제목 s2v, 실험 이름 clean\_street2\_video), Track B \texttt{track-b-4} (제목 drive-dotonbori, 실험 이름 b6).

% ---------------------------------------------------------------------------
\section*{들어가며: 이 리포트를 읽는 법}

이 과제에서 나는 두 트랙을 한 줄의 이야기로 진행했다. Track A에서는 정해진 모델(ResNet18, 224×224)과 정해진 사진 안에서, 사진을 고르고 빼거나 학습 방법을 바꾸는 것만 할 수 있었다. 그 과정에서 모델이 왜 틀리는지를 계속 재고 보았고, 결국 ``사진을 빼거나 가리는 것으로는 부족하고, 학습에 없던 장면의 영상 자체가 필요하다''는 결론에 이르렀다. Track B에서는 그 결론을 직접 시험했다. 더 큰 사진으로 글자를 보게 하고, 부족한 장면의 유튜브 영상을 직접 골라 더했다. 그 과정에서 모델이 엉뚱한 지름길을 배우는 것을 발견해 고쳤고, 마지막에는 오답을 보고 부족한 장면을 정확히 채워, 새 검증에서 가장 높은 점수를 얻었다.

\paragraph{이 리포트에서 쓰는 점수.} 모든 점수는 도시마다 정답률을 구한 뒤 다섯 도시를 평균한 값이다.
\begin{itemize}
\item \textbf{공식 검증}: 스타터 검증 사진 12,077장의 점수. 이 사진의 약 70\%는 위키미디어 공용(commons)의 사진이라, 명소나 건물을 골라 찍은 사진이 많다.
\item \textbf{영상 프레임 점수}: 공식 검증 사진 중 유튜브 영상의 프레임만 잰 점수. 검증 영상은 학습 영상과 겹치지 않으므로 ``처음 보는 영상을 얼마나 맞히는가''를 뜻한다. 리더보드 사진도 모두 유튜브 걷기·운전 영상의 프레임이라, 리더보드와 같은 종류다.
\item \textbf{학습 프레임 점수}: 학습에 쓴 유튜브 프레임을 다시 맞힌 비율. 이것과 영상 프레임 점수의 차이를 \textbf{외우기}의 크기로 썼다.
\item \textbf{새 검증}: Track B를 위해 직접 만든 검증 세트(도시마다 140장). 스타터 데이터나 내 학습 영상과 채널이 겹치지 않는 유튜브 영상에서 뽑았다.
\item \textbf{리더보드}: Track A는 채점하는 쪽이 내 레시피를 seed 0, 1, 2로 세 번 학습한 평균과 범위, Track B는 내가 올린 모델 그대로의 점수다.
\end{itemize}

모든 숫자는 맨 끝의 \texttt{experiments.csv}와 \texttt{LOG.md}, 그리고 \texttt{outputs/}의 출력 파일에서 가져왔다. 각 실험의 이유와 해석은 LOG에 그때그때 적었고, 이 리포트는 그것을 하나의 흐름으로 엮은 것이다.

% ---------------------------------------------------------------------------
\section*{Track A}

\begin{center}\footnotesize
\begin{tabular}{lp{0.6\linewidth}r}
\toprule
실험 & 바로 앞 실험에서 바뀐 것 & 공식 검증 \\
\midrule
baseline & 시작 레시피 그대로 (seed 0 / seed 1) & 0.489 / 0.496 \\
balanced & 다섯 도시를 같은 비율로 뽑기 & 0.514 \\
balanced\_no\_flip & balanced에서 좌우 반전 끄기 & 0.506 \\
balanced\_crop60 & balanced에서 잘라 내기 최소 넓이 35\%→60\% & 0.498 \\
balanced\_clean & balanced에 사진 정제 1차(1,074장 제외) & 0.521 \\
balanced\_clean\_mixup & balanced\_clean에 mixup & 0.506 \\
clean2 & balanced\_clean에 정제 2차(동물·자동차·음식 854장 더 제외) & 0.510 \\
clean2\_food & clean2에서 음식 306장을 다시 넣음 & 0.512 \\
clean\_street2 & balanced\_clean에서 거리 사진 2배 & 0.508 \\
clean\_street3 & 거리 사진 3배 & 0.494 \\
clean\_street2\_video & 거리 사진 2배 + 영상마다 똑같이 뽑기 & 0.514 \\
clean\_street2\_video\_lr & 위에 층별 학습률 & 0.498 \\
clean\_street2\_video\_color & 위에 색 변화 강화 (릴리스 weight-color) & 0.516 \\
clean\_color & balanced\_clean에 색 변화 강화만 & 0.510 \\
clean3\_s2v\_color & weight-color에 학습 거리 사진 정제 3차(834장 제외) & 0.511 \\
clean3\_blur & clean3\_s2v\_color에 흐림 & 0.494 \\
clean3\_cutmix & clean3\_s2v\_color에 CutMix & 0.494 \\
\bottomrule
\end{tabular}
\end{center}

\subsection*{1. 출발점: 모델은 사진이 많은 도시로 쏠린다}

먼저 시작 레시피를 seed만 바꿔 두 번 학습했다. 전체 점수는 0.489와 0.496으로 거의 같아, 이후 모든 비교에서 약 0.6\%p를 ``우연히 생길 수 있는 차이''의 기준으로 썼다. 도시별로 보면 점수 순서가 학습 사진 수와 거의 같았다. 사진이 가장 많은 도쿄와 서울이 가장 높았고, 가장 적은 부산은 낮았다. 특히 오사카와 타이베이는 사진 수가 비슷한데도 0.33\textasciitilde 0.38과 0.53으로 크게 달랐다. 여기서 첫 가설을 세웠다. 모델은 어느 나라인지는 대강 맞히지만, 헷갈리는 사진은 같은 나라에서 사진이 많은 도시로 보낸다. 그래서 오사카는 도쿄에, 부산은 서울에 밀리고, 같은 나라에 경쟁 도시가 없는 타이베이는 그런 손해를 보지 않는다.

이를 확인하려고 다섯 도시를 학습에서 같은 비율로 뽑았다(balanced). 부산은 8.6\%p, 오사카는 8.8\%p 올라 가설과 맞았다. 그런데 예상과 달리 타이베이도 8.4\%p 올랐다. 쏠림은 같은 나라 안에서만이 아니라, 나라와 관계없이 사진이 많은 도시 쪽으로 일어나고 있었던 것이다. 다만 전체 점수는 2.1\%p만 올랐다. 대부분은 서울과 도쿄의 점수가 다른 도시로 옮겨 간 것이고, 도시를 구분하는 능력 자체가 나아진 것은 아니었다. 이후 모든 Track A 실험은 이 도시 균등을 기본으로 했다.

\subsection*{2. 글자와 잘라 내기: 224 사진의 한계를 처음 만나다}

baseline의 오답에는 간판의 한글이나 가나가 분명히 보이는데도 다른 나라로 틀린 사진이 많았다. 학습 때 사진을 무작위로 좌우 반전하면 글자와 좌측통행이 뒤집히므로, 이를 끄면 모델이 글자와 통행 방향을 더 잘 쓸 것이라고 생각했다(no\_flip). 그러나 점수는 0.8\%p 내려갔고, 일본어와 좌측통행이 핵심 단서인 도쿄와 오사카도 그대로였다. 간판이 보이는 오답 20장을 모델이 실제로 받는 224×224로 줄여 직접 보니, 글자가 뭉개져 있었다. 반전 때문에 글자를 못 쓴 것이 아니라, 처음부터 글자를 읽을 수 없었던 것이다. 이미지 크기가 고정된 Track A에서는 이를 고칠 수 없으므로, 글자는 Track B에서 다루기로 했다.

학습 때 사진을 넓이의 35\%까지 잘라 확대하면, 작은 조각에는 산, 바다, 전선처럼 사진 전체를 봐야 보이는 단서가 빠진다. 그래서 최소 넓이를 60\%로 올렸다(crop60). 그런데 학습 loss는 0.782에서 0.662로 줄었는데 검증은 1.6\%p 내려갔다. 사진이 매번 비슷하게 보여 외우기 쉬워진 것으로 보았다. 이것이 ``학습은 잘 되는데 처음 보는 사진은 못 맞힌다''는 신호를 처음 본 실험이었다.

\paragraph{오답이 그럴듯해진 것은 착각이었다.} balanced의 오답을 보면서 ``도시의 특징을 더 잘 파악한 상태에서 틀린다''는 느낌을 받았다. 이를 확인하려고 정답이 2순위 안에 드는 비율(top-2)을 쟀더니, baseline 0.721, balanced 0.728로 거의 같았다. 쏠림이 줄면서 오답 모음의 구성이 바뀌어 생긴 착각이었다. 대신 새로 알게 된 것이 있다. balanced에서 틀린 사진의 약 44\%는 정답을 2순위로 두고 있었다. 모델은 도시 단서를 꽤 잡고 있지만, 1순위와 2순위 사이에서 확신하지 못하는 경우가 많았다.

\subsection*{3. 사진 정제: 단서 없는 사진을 빼면 나아질까}

오답을 직접 보며 규칙을 세우고 학습 사진을 세 번 정제했다. 자세한 과정은 다음 절(``사진을 보고 내린 결정'')에 있다. 1차(1,074장)는 공식 검증을 0.514에서 0.521로 올렸지만, seed 차이와 비슷한 크기였다. 정제 후 남은 학습 사진을 보니 차, 동물, 접시 같은 노이즈가 여전히 많아, 2차에서는 사진 내용으로 후보를 찾아 854장을 더 뺐지만 0.510으로 오히려 조금 내려갔다. 리더보드도 clean 0.415, clean2 0.395로 오르지 않았다.

정제가 효과를 보이지 않자 학습 과정을 의심했다. 모델이 학습 사진은 87\%를 맞히는데 검증 사진은 51\%만 맞혀, 차이가 36\%p였다. 덜 학습된 것이 아니라 학습 사진을 외운 것으로 보고, 사진 두 장을 겹쳐 외우기 어렵게 하는 mixup을 넣었다. 학습 정답률은 0.883에서 0.848로 내려가 외우기는 실제로 줄었지만, 검증은 1.5\%p 내려갔다. 이 결과에서 처음으로 ``외우기를 줄이는 것만으로는 점수가 오르지 않는다''는 단서를 얻었다.

\subsection*{4. 거리 사진에 집중하기: 문제는 같은 영상의 반복이었다}

clean과 clean2의 리더보드를 보니, 학습 사진 중 거리 사진(유튜브 프레임, osv5m)의 비율이 높은 도시일수록 리더보드 점수가 높았다(거리 사진 비율 도쿄 48.3\%, 서울 41.1\% / 리더보드 .63, .53, 나머지 세 도시는 22.9\textasciitilde 33.1\% / .26\textasciitilde .29). 리더보드는 거리 영상이므로, 거리 사진을 더 자주 보면 점수가 오를 것으로 예상했다.

그런데 거리 사진을 2배, 3배 자주 뽑자 리더보드와 같은 종류인 영상 프레임 점수가 0.418에서 0.376, 0.358로 오히려 내려갔다. 흥미로운 것은 같은 거리 사진인 osv5m 사진의 점수는 2배에서 2.0\%p 올랐다는 점이다. 유튜브 프레임은 소수의 영상에서 대량으로 나오지만 osv5m 사진은 한 장씩 따로 찍혔다. 그래서 거리 사진 자체가 나쁜 것이 아니라, \textbf{같은 영상을 반복해서 보는 것}이 문제라고 보았다. 영상마다 똑같이 뽑자 프레임 쏠림이 가장 컸던 부산과 타이베이가 5.9\%p씩 회복되어, 이 생각과 맞았다.

이를 직접 확인하려고 학습 프레임과 처음 보는 영상의 프레임을 따로 쟀다. 모델은 학습한 영상의 프레임을 97.1\% 맞히면서 처음 보는 영상은 38.8\%만 맞혔다. 부산은 학습 프레임을 98.6\% 맞히면서 처음 보는 영상은 18.2\%만 맞혔고, 처음 보는 부산 영상의 45.0\%를 서울로 답했다. 모델은 도시가 아니라 영상 하나의 모습(카메라, 날씨, 색감)을 외우고 있었고, 처음 보는 거리에서는 그 기준이 통하지 않아 영상이 많은 같은 나라의 도시로 답했다. 도시별 유튜브 영상 수(도쿄 36, 서울 34, 오사카 18, 부산 17, 타이베이 9)를 보니, 리더보드에서 .6을 넘은 도시는 영상이 34개 이상인 두 도시뿐이었다. 점수를 가른 것은 거리 사진의 비율이 아니라 \textbf{서로 다른 영상의 수}일 수 있었다.

\subsection*{5. 외우기 줄이기: 효과는 조건에 따라 달랐다}

Track A에서는 영상을 더할 수 없으므로, 외우기를 줄이는 방법을 시험했다. 층별 학습률은 모든 점수를 내렸다(공식 검증 −1.6\%p). 색 변화 강화는 가중치 모델에서 영상 프레임 점수를 2.4\%p 되찾고 외우기 차이도 0.583에서 0.548로 줄였다. 그래서 이 모델(weight-color)을 릴리스했다. 그런데 같은 색 변화를 가중치가 없는 모델에 넣자 영상 프레임 점수가 2.5\%p 내려갔다. 색 변화는 그 자체로 좋은 방법이 아니라, 외우기가 심할 때만 도움이 되는 방법이었다.

리더보드에서는 weight-color가 0.410으로, 직전 릴리스 s2v(0.418)와 구별되지 않았다. 여기서 ``우리 검증으로 리더보드를 예측할 수 있는가''를 따로 쟀다. 릴리스 3개의 영상 프레임 점수와 리더보드를 비교하면, 어느 도시가 어려운지는 잘 맞혔지만(도시별 점수 15개, r = 0.81), 모델을 바꿨을 때 도시별로 오르고 내리는 방향은 전혀 맞히지 못했다(r = −0.11). 레시피 변경이 만드는 1\textasciitilde 3\%p 차이는 채점의 흔들림만큼 작아서, 어떤 검증으로도 순서를 가리기 어려웠다.

\paragraph{검증이 어긋나는 방향은 늘 같았다.} 도시마다 검증 점수와 리더보드의 차이를 보니, 부산과 도쿄는 늘 검증이 낮고 타이베이와 오사카는 늘 높았다. 가장 많이 어긋나는 타이베이와 부산의 학습 거리 사진과 검증 거리 사진을 직접 비교했다. 타이베이 학습 영상은 거의 낮의 일상 브이로그였지만, 검증 영상은 밤의 야시장, 네온 거리, 야간 운전이었다. 부산 학습 영상은 해운대 해변, 해안 산책로, 공원이 많았지만, 검증 영상은 상가, 시장, 도심 도로였다. 도시마다 검증 영상이 4\textasciitilde 6개뿐이라, 그 몇 개가 어떤 장면이냐에 따라 점수가 크게 달라졌다. 그리고 학습 거리 사진을 영상별로 보다가, 셀카, 실내, 음식 같은 노이즈가 많다는 것도 발견했다. 거리 사진을 2배로 뽑고 있으니 이 노이즈도 2배로 학습되고 있었다.

\subsection*{6. 마지막 시도와 결론}

그래서 다섯 도시의 학습 거리 사진을 모두 직접 보고 834장을 뺐다(clean3, 다음 절). 학습 프레임은 여전히 96\%를 맞혔고, 외우기 차이는 0.548에서 0.560으로 줄지 않았다. 모델이 외우는 것은 셀카 같은 특정 사진이 아니라, 한 영상의 모든 프레임에 남는 선명도, 색감, 날씨 같은 ``영상의 모습''인 것으로 예상했다. 도시별로는 서울(−8.1\%p)과 부산(+6.6\%p)이 정반대로 움직여 한국 안의 서울 쏠림이 줄었고, 일본에서는 도쿄 사진을 오사카로(.178→.212), 오사카 사진을 도쿄로(.333→.352) 답하는 비율이 함께 늘어 두 도시를 가르는 힘이 약해졌다. 전체 차이는 흔들림 범위였지만, 리더보드가 낮은 부산과 타이베이의 영상 프레임 점수가 Track A 모델 중 가장 높아서 이 모델을 \texttt{track-a-7}(clean3)로 릴리스했다. 그러나 이 릴리스는 리더보드에서 채점되지 않아, 최종 릴리스는 s2v로 바꿨다(아래).

영상의 모습을 가리려고 흐림과 CutMix도 시험했다. CutMix는 학습 프레임 정답률을 .962에서 .918로 낮춰 외우기를 3.6\%p 줄였지만, 영상 프레임 점수는 오르지 않았다(−0.8\%p). 흐림은 외우기도 줄이지 못하고 점수만 2.6\%p 잃었다. 흐림이 영상의 선명도뿐 아니라 간판과 건물의 세부 단서까지 지운 것으로 예상했다.

이 실험들을 모아 Track A의 결론을 내렸다.
\begin{itemize}
\item Track A 모델은 224 사진에서 글자를 읽지 못한다. 기댈 수 있는 단서는 간판의 모양과 색, 도로와 차선, 건물 같은 장면의 특징뿐이다.
\item 그런데 학습 영상의 종류가 적어서, 모델은 이런 나라별 특징보다 각 영상에만 있는 특징을 먼저 배운다. 예를 들어 운전 영상이 몇 개뿐이면, 모델은 도로의 모습보다 차 안의 모습이나 하늘로 영상을 구분한다. 학습한 운전 프레임을 거의 100\% 맞힌 것도 이와 맞는다.
\item 사진을 빼거나(clean3) 가리는 것(색 변화, 흐림, CutMix)으로는 이를 막지 못했다. 처음 보는 장면을 알아보려면 비슷한 장면을 학습에서 본 적이 있어야 한다.
\item 학습 사진이 다양하지 않은 도시는 처음 보는 장소에서 밀린다. 모델은 처음 보는 부산 영상 프레임의 약 절반을 서울로, 오사카 영상 프레임의 약 3분의 1을 도쿄로 답했다. 이것은 레시피로 고칠 수 없는 데이터 부족의 문제라고 판단했다.
\end{itemize}
Track A에서는 사진을 더할 수 없으므로, 이 결론은 Track B에서 시험했다.

\paragraph{최종 레시피와 리더보드.} clean3를 담은 \texttt{track-a-7}은 리더보드에서 ``no recipe.py in the release''로 채점되지 않았다. 릴리스를 레시피 파일 없이 먼저 만들었고, 나중에 \texttt{recipe.py}와 \texttt{selection.csv}를 붙였지만 채점기는 이 릴리스를 다시 채점하지 않았다. clean3는 리더보드 점수를 한 번도 받지 못해, 최종 사진에서 스타터보다 낮게 나와도 알 수 없다. 적어도 리더보드에서 확인된 레시피를 최종으로 올려야 한다고 판단해, \texttt{track-a-5}(s2v)와 같은 파일로 \texttt{track-a-8}(s2v)을 다시 릴리스했다. s2v는 우리 Track A 릴리스 중 리더보드가 가장 높았고(0.418), 스타터 레시피를 그대로 쓴 다른 학생의 점수(0.407)보다도 높다.

최종 레시피 s2v: 다섯 도시 균등, 거리 사진 2배, 영상마다 똑같이 뽑기, 그리고 정제 1차를 거친 사진 43,183장.

\begin{center}\footnotesize
\begin{tabular}{lrrrrrrr}
\toprule
릴리스 (실험) & 공식 검증 & 리더보드 & 서울 & 부산 & 도쿄 & 오사카 & 타이베이 \\
\midrule
clean (balanced\_clean) & 0.521 & 0.415 & .61 & .25 & .65 & .29 & .27 \\
clean2 (clean2) & 0.510 & 0.395 & .53 & .27 & .63 & .29 & .26 \\
s2v (clean\_street2\_video) & 0.514 & 0.418 & .59 & .27 & .65 & .28 & .30 \\
weight-color (clean\_street2\_video\_color) & 0.516 & 0.410 & .54 & .23 & .66 & .33 & .28 \\
clean3 (clean3\_s2v\_color) & 0.511 & \multicolumn{6}{l}{채점되지 않음 (no recipe.py in the release)} \\
\bottomrule
\end{tabular}
\end{center}

Track A의 모든 릴리스는 0.395\textasciitilde 0.418 안에 있다. 같은 레시피를 seed만 바꿔도 clean은 0.399\textasciitilde 0.444로 4.5\%p 움직였으므로, 어떤 변경도 채점의 흔들림과 구별되는 효과를 보이지 못했다. 공식 검증은 리더보드보다 늘 10\%p 이상 높다. 공식 검증의 약 70\%가 위키미디어 사진이고 리더보드는 모두 거리 영상이기 때문인 것으로 예상하며, 리더보드와 같은 종류인 영상 프레임 점수는 clean에서 0.418로 리더보드(0.415)와 거의 같았다.

% ---------------------------------------------------------------------------
\section*{사진을 보고 내린 결정}

\subsection*{도시를 가르는 단서와 오답의 공통점}

처음에 도시마다 학습 사진 100장을 보고 도시를 가르는 단서를 적었다. 서울은 고층건물, 회색, 흐린 하늘, 한글 간판, 한옥, 빽빽한 건물, 건물 뒤편 산. 부산은 바다, 한글 간판, 파란 버스, 산비탈에 층층이 들어선 주택. 도쿄는 일본 간판, 좌측통행 차도, 촘촘한 전봇대와 전선. 오사카는 일본 간판, 도톤보리 강, 밤에 빛나는 간판 거리, 낮은 건물. 타이베이는 낮은 건물, 야시장, 많은 오토바이와 스쿠터, 울창한 아열대 나무.

baseline이 틀린 사진을 다섯 도시 모두 보니 두 가지가 반복됐다(그림~\ref{fig:mistakes}). 하나는 도시를 알 수 없는 사진이다. 차나 음식 클로즈업, 배경이 없는 실내, 둥글게 휜 360도 파노라마, 차 안에서 찍은 블랙박스 화면, 불빛만 보이는 밤 사진이었다. 모델은 이런 사진에서 건물이나 거리보다 촬영 조건을 보고 도시를 고르는 것 같았다. 다른 하나는 간판의 글자가 보이는데도 다른 나라로 틀린 사진이다.

\begin{figure}[h]\centering
\picturefile[0.8\linewidth]{figures/mistakes_baseline_seed0_busan_00.png}
\caption{baseline이 틀린 부산 검증 사진. 클로즈업, 실내, 블랙박스 화면처럼 도시를 알 수 없는 사진이 반복된다.}
\label{fig:mistakes}
\end{figure}

\subsection*{행동하기 전에 쓴 규칙}

규칙: 도시를 알아볼 단서가 없는 학습 사진을 뺀다.
\begin{enumerate}
\item 도시의 거리나 장소가 아닌 클로즈업: 물건, 음식, 인물, 지도 등이 화면 대부분을 차지하는 사진. 거리가 잘 보이면 남긴다. (2차에서 동물을 더했다.)
\item 360도 파노라마: 도시 모습을 왜곡해, 촬영 방식 자체를 외울 위험이 있다.
\item 무엇을 찍었는지 알 수 없을 만큼 어둡거나 흐린 사진. 어두워도 간판이나 야시장이 보이면 남긴다.
\item 코스프레 사진: 후보를 보다가 더했다. 모델은 도쿄 학습 사진 전체 중 6.3\%만 타이베이로 답했지만, 도쿄 코스프레 사진은 39\%(46장 중 18장)를 타이베이로 답했다(그림~\ref{fig:cosplay}). 타이베이 학습 사진에 코스프레 사진이 많아, 코스프레 복장 자체를 타이베이의 단서로 배운 것으로 보았다.
\end{enumerate}
예외: 그 나라 글자가 명확히 보이면 남긴다.

\begin{figure}[h]\centering
\picturefile[0.8\linewidth]{figures/cosplay_said_taipei_00.png}
\caption{타이베이가 아닌데 모델이 타이베이로 답한 코스프레 학습 사진. 이를 보고 규칙 4를 더했다.}
\label{fig:cosplay}
\end{figure}

\subsection*{세 번의 정제와 그 결과}

\textbf{1차 (balanced\_clean).} 사진 제목에 음식, 박물관, 코스프레 같은 말이 있는 사진, 파노라마 비율의 사진, 가장 어두운 사진, 모델이 가장 의심한 사진을 후보로 모아 직접 보고 1,074장을 뺐다(클로즈업 686, 파노라마 162, 어두운 사진 46, 코스프레 180). 공식 검증은 0.514에서 0.521로 올랐지만 seed 차이와 비슷했다. 코스프레를 뺐는데도 도쿄 사진을 타이베이로 답한 수는 449장에서 471장으로 줄지 않았다. 검증 세트의 타이베이 폴더에도 코스프레 사진이 있어서, 코스프레는 학습과 검증 모두에서 타이베이와 묶여 있었다. 그래도 도시의 특징이 아닌 잘못된 단서이므로 빼는 것이 맞다고 판단했다.

\textbf{2차 (clean2).} 제목으로 찾는 방법은 제목에 그런 말이 없는 노이즈를 놓쳤다. 그래서 ImageNet 분류기(학습 없이 그대로 사용)로 1순위가 음식, 동물, 자동차인 사진을 후보로 모으고, 모두 직접 보고 854장을 뺐다. 음식은 나라별 특징일 수도 있어서 효과가 모호하다고 보고, 음식만 다시 넣은 clean2\_food도 따로 학습해 효과를 떼어 보았다. 세 가지를 모두 뺀 효과는 −1.1\%p, 음식만의 효과는 −0.2\%p로 모두 흔들림 범위였다.

\textbf{3차 (clean3\_s2v\_color).} 리더보드와 같은 종류인 학습 거리 사진을 다섯 도시 모두 영상 순서대로 직접 보고, 834장을 뺐다(그림~\ref{fig:selfie}). 같은 구도가 반복되는 셀카, 버스 안에서 찍은 밤 영상, 비정상적으로 파랗게 찍힌 영상, 자동차 전시장과 관광지 실내였다. 결과는 Track A 6번에 적었다.

\begin{figure}[h]\centering
\picturefile[0.8\linewidth]{figures/look_street_taipei_12.png}
\caption{타이베이 학습 거리 사진 시트. 같은 구도의 셀카가 반복되어(12--22번, 29--36번) 뺐다.}
\label{fig:selfie}
\end{figure}

최종 릴리스 s2v의 선택은 다음과 같다(\texttt{curate.py count}의 출력). 정제 1차(1,074장)까지만 반영되어 있고, 2차와 3차에서 뺀 사진은 clean2와 clean3 실험에만 쓰였다(Track A 6절).
\outputfile{outputs/curate_count.txt}

세 번 모두 점수는 흔들림 범위 안에서 움직였다. 뺀 사진이 매번 전체의 2\% 안팎으로 적었고, 검증 사진에도 같은 종류의 노이즈가 남아 있었으며, 모델이 외우는 것이 특정 사진이 아니라 영상의 모습이었기 때문인 것으로 예상한다. 그래도 정제 기준은 Track B에서 영상을 모을 때 그대로 썼다.

\subsection*{에이전트에게 한 요청 (대화에서 그대로 옮김)}
\begin{itemize}
\item ``4plW2VUTBi8 전체 빼줘 거의 다 노이즈야. 인물 얼굴 사진 혹은 메뉴판 같은 사물 사진이니까. 특히 인물 사진의 경우 계속 비슷한 셀카 구도이기 때문에 검증에서 비슷한 구도의 셀카 이미지가 나오면 타이베이로 잘못 분류할 가능성도 있어.''
\item ``cjO9fH2tRyk (버스 안) 이거 전체 삭제할게. 버스 안이라서 버스 창의 도시 배경보다 버스 내부에서 찍은 프레임 자체를 외울 위험이 있고 버스 창으로 보이는 것도 밤이라서 어두운 도로 뿐이라 부산 도시에 대한 단서는 보이지 않아.''
\item ``Vi7ozoEThsc 전체 삭제. 이미지가 모두 비정상적으로 파랗게 찍혔어. 이는 색으로 나라를 분류하는 오학습을 유도할 수 있기에 삭제.''
\item ``mepi1nhmiom은 다 지워. 360도 이미지야. 타이페이 03의 13\textasciitilde 24 너무 어두우니까 지워.''
\item ``도톤보리 낮영상만 넣을거야. 밤은 오답 거의 없어.''
\item ``3도시 각각 보완 영상 찾아줘. 이때 영상 학습이 너무 많이 되면 또 도쿄와 서울이 뒤집힐 수 있어. 각 도시의 영상 학습 채널 종류와 정도가 비슷할수록 좋을듯.''
\end{itemize}

% ---------------------------------------------------------------------------
\section*{Track B}

\begin{center}\footnotesize
\begin{tabular}{llrrr}
\toprule
모델 & 바로 앞 모델에서 바뀐 것 & 공식 검증 & 새 검증 & 리더보드 \\
\midrule
b1 & 시작 코드 그대로 (ResNet18, 224, 4 epoch) & 0.493 & & \\
b2 & MobileNetV3-Large, 448, 도시 균등, 6 epoch & 0.585 & 0.474 & 0.432 \\
b3 & 부산·오사카·타이베이에 추가 영상 프레임(720p) & 0.574 & 0.491 & \\
b4 & 추가 프레임을 스타터와 같은 480p로 & 0.569 & 0.456 & \\
b5 & 서울·도쿄에도 같은 종류의 영상 추가 & 0.580 & 0.486 & 0.468 \\
b6 & 운전 영상과 도톤보리 낮 영상 추가 & 0.579 & 0.523 & 0.478 \\
b7 & b6에서 거리 사진 2배 & 0.582 & 0.513 & \\
\bottomrule
\end{tabular}
\end{center}
새 검증 점수는 모두 480p로 다시 저장한 새 검증으로 잰 값이다(b3 참고).

\subsection*{1. 큰 사진으로 글자를 보게 하기 (b1, b2)}

Track A에서 224 사진의 간판 글자는 뭉개져 있었다. Track B는 파라미터 1,200만 개, 연산 20억 번만 넘지 않으면 무엇이든 바꿀 수 있으므로, 작은 네트워크(MobileNetV3-Large)로 큰 사진(448)을 보게 했다(연산 0.85B). Track A에서 효과가 확인된 도시 균등도 넣었다. 공식 검증은 시작 코드(b1, 0.493)보다 9.2\%p 오른 0.585로, 모든 Track A 결과를 넘었다. 가장 많이 오른 도시는 Track A에서 가장 낮던 부산(+15.2\%p)과 오사카(+16.9\%p)였다. 다만 세 가지를 한 번에 바꿔 무엇이 얼마나 기여했는지는 가를 수 없다.

\paragraph{b2는 정말 글자를 볼까.} b2의 오답을 보니 ``한자가 있으면 헷갈려하고 한글이 있으면 대부분 한국으로'' 답했다. 이를 두 가지로 확인했다. 첫째, 검증 사진의 글자를 OCR로 찾아 종류별로 나누니, 글자가 많은 사진에서 나라를 맞힌 비율이 한글 0.883, 가나 0.821로 글자 없는 사진(0.768)보다 높았다. 둘째, 글자가 원래 쉬운 사진에만 있을 수도 있어서, 같은 사진에서 글자만 회색으로 가려 보았다. 같은 넓이의 다른 곳을 가렸을 때는 4.5\%가 맞다가 틀렸지만, 글자를 가리면 9.9\%가 틀렸다(z = 5.0). b2는 글자를 단서로 쓴다. 그러나 글자를 가려도 전체 정답률은 1\%p만 떨어져, 대부분의 사진은 글자 없이 판단했다. 오답을 직접 본 내 결론은 ``내가 봐도 구분이 쉽지는 않아. 언어 구분을 할 수 있지 않는 한 말이지''였다. 글자가 없는 장면에서도 도시를 가르려면, 다양한 장면의 학습 사진이 더 필요했다.

\subsection*{2. 리더보드를 닮은 검증 만들기}

b2(format)의 리더보드는 0.432였는데, 공식 검증은 0.585였다. 앞에서 보았듯 공식 검증은 리더보드와 사진 종류가 달랐고, 영상 프레임 점수도 도시마다 검증 영상이 4\textasciitilde 12개뿐이라 비슷한 모델 사이의 순서를 가리지 못했다. 예를 들어 format의 오사카는 영상 프레임 점수 0.530에 리더보드 0.30이었다. 또 리더보드 사진은 걷기 영상과 운전 영상에서 나오는데, 스타터 영상을 제목으로 나눠 보니 오사카와 타이베이에는 운전 영상이 학습과 검증 모두에 하나도 없었다. 그래서 리더보드의 조건을 닮은 \textbf{새 검증}을 만들었다. 다섯 도시마다 걷기 영상 4\textasciitilde 5개와 운전 영상 3개를 찾고, 스타터 데이터와도 학습용 영상과도 채널이 겹치지 않게 했으며, 리더보드처럼 영상당 최대 20장씩 도시마다 140장을 뽑았다.

\subsection*{3. 추가 영상을 모은 방법}

Track A가 가리킨 부족한 점, 즉 영상 수, 학습에 없던 장면, 운전 영상을 채우려고 유튜브 영상 링크를 직접 찾았다.
\begin{itemize}
\item 도시마다 밤과 낮, 여러 지역, 비 오는 날, 운전 영상이 들어가도록 골랐다. 스타터 데이터에 있는 영상과, 검증 영상과 채널이 겹치는 영상은 뺐다.
\item 영상마다 1분에 한 장씩 프레임을 뽑고, 인트로와 엔딩이 있는 앞뒤 5\%는 뺐다. 다른 영상보다 훨씬 긴 영상은 60장으로 맞췄다.
\item 모든 프레임은 스타터와 같은 854×480, JPEG 품질 85로 저장했다(이유는 아래 b3).
\item Track A에서 쓴 기준으로 시트를 직접 보고, 사물·인물·실내 프레임과 문제가 있는 영상을 뺐다(그림~\ref{fig:dotonbori}). 비 오는 날 렌즈에 빗물이 묻어 도시가 안 보이는 영상, 카메라 필터로 색감이 다른 영상, 360도 영상, 화면 위에 자막과 국기가 계속 나오는 영상이었다. 이유는 \texttt{scripts/mine\_dropped.csv}와 \texttt{scripts/mine\_removed\_videos.csv}에 적었다.
\end{itemize}

\begin{figure}[h]\centering
\picturefile[0.8\linewidth]{figures/mine_osaka_dotonbori_00.png}
\caption{b6를 위해 넣은 도톤보리 낮 영상의 프레임. 이 시트에서 실내·사물 프레임을 골라 뺐다.}
\label{fig:dotonbori}
\end{figure}

\subsection*{4. 영상을 넣자 지름길이 생기다 (b3, b4)}

먼저 영상이 적은 부산(12개), 오사카(11개), 타이베이(11개)에만 영상을 넣었다(b3). 그러자 새 검증(이때는 720p로 뽑았다)에서 서울은 .607에서 .029로, 도쿄는 .793에서 .071로 무너졌고, 서울 사진은 부산으로, 도쿄 사진은 오사카로 정확히 같은 나라 안에서 옮겨 갔다. 공식 검증은 1\%p만 내려가서, 새 검증이 없었다면 알아채지 못했을 것이다.

원인은 화질로 예상했다. 스타터 데이터의 영상 프레임은 854×480에 JPEG 품질 약 85로 압축되어 있지만, 내가 뽑은 프레임은 720p로 훨씬 선명했다. 추가 프레임은 세 도시에만 넣었으므로 학습에서 선명한 영상 프레임은 언제나 이 세 도시였고, 모델은 ``선명하면 부산·오사카·타이베이''를 배운 것으로 보였다. 이는 Track A에서 본 ``영상의 모습을 외우는'' 현상과 같다. 휴대폰 사진도 선명하므로 최종 순위에서 특히 위험했다. 그래서 모든 추가 프레임과 새 검증 사진을 스타터와 같은 480p, 품질 85로 다시 저장했다. 모델은 짧은 변을 448로 줄여 보므로, 글자를 읽는 데 필요한 해상도는 거의 잃지 않는다.

검증 사진만 480p로 바꿔도 b3의 서울이 .029에서 .321로, 도쿄가 .071에서 .529로 회복되어, 화질은 지름길의 일부였다. 그러나 학습 프레임까지 480p로 맞춘 b4에서도 서울은 .150, 도쿄는 .393에 머물렀고, b4는 서울 사진의 62\%를 부산으로 답했다. 남은 원인은 영상의 출처인 것으로 예상했다. 추가 영상은 모두 최근의 걷기 투어 채널에서 나온 비슷한 촬영 스타일이었고, 이런 영상이 세 도시에만 있었다. 이를 고치려면 서울과 도쿄에도 같은 종류의 영상을 넣어, 그 스타일이 다섯 도시 모두에 있게 해야 했다.

\subsection*{5. 다섯 도시를 같은 조건으로 (b5)}

서울(8개)과 도쿄(10개)에도 같은 종류의 영상을 넣었다. 서울과 도쿄는 인프라가 발달해 장소가 다양하고 배울 장소도 많다고 보았고, 영상마다 똑같이 뽑기는 하지 않되 서울과 도쿄만 너무 높아지면 그때 균형을 맞추기로 했다. 새 검증에서 서울은 .150에서 .464로, 도쿄는 .393에서 .614로 회복됐고, 서울 사진을 부산으로 답한 비율도 .621에서 .357로 줄었다. 리더보드는 0.432에서 0.468로 올랐고, 다섯 도시가 모두 올랐다. Track A의 어떤 변경도 만들지 못한, 흔들림보다 큰 상승이었다.

그런데 새 검증에서는 서울과 도쿄가 b2보다 낮았는데 리더보드에서는 둘 다 올랐다. Track A에서 본 것처럼, 모델을 바꿨을 때의 도시별 변화는 검증으로 예측되지 않았다. 타이베이는 새 검증 .600인데 리더보드 .22로 여전히 가장 낮았다.

\subsection*{6. 오답을 보고 부족한 장면을 채우기 (b6)}

b5의 오답을 직접 보니, 한자가 잘 보이는데도 타이베이 사진을 서울로 답한 경우가 많았다. 모델이 글자를 충분히 쓰지 못하는 것인지 확인하려고, b2에 했던 글자 확인을 b5로 다시 했다(\texttt{outputs/text\_check\_b2\_b5.txt}). b5는 글자를 가리면 11.6\%가 맞다가 틀려 다른 곳을 가렸을 때(4.7\%)의 2.5배였고, b2(9.9\% 대 4.5\%)보다 글자에 조금 더 기댔다. 글자가 많은 타이베이 한자 사진의 정답률도 0.780에서 0.854로 올라, 서울로 답한 비율이 .146에서 .049로 줄었다. 그러니 b5는 큰 간판의 한자는 구분했다. 새 검증의 오답은 간판이 작거나 멀거나 가운데 정사각형 밖에 있어 글자 대신 거리 모습으로 판단한 사진으로 보였다. 448은 리더보드 사진의 원래 크기(짧은 변 480)를 거의 그대로 보는 크기라, 이미지를 더 키우는 것은 답이 아닐 것으로 판단했다.

그래서 b5의 새 검증 오답을 영상별로 나눠 보았다. 리더보드가 낮은 세 도시 모두, 운전 영상의 프레임을 걷기 영상보다 훨씬 많이 틀렸다(틀린 비율: 부산 67\% 대 36\%, 오사카 88\% 대 69\%, 타이베이 63\% 대 23\%. 정답률로는 표~
\ref{tab:walkdrive}의 b5 값). 특히 도심 도로 운전 영상은 85\textasciitilde 90\%를 틀렸고, 같은 나라에서 학습 영상이 많은 서울과 도쿄로 갔다. Track A에서 예상한 대로, 운전 영상이 몇 개뿐이면 도로의 특징보다 영상마다의 모습을 배워 처음 보는 도로에서 틀리는 것으로 예상된다. 학습 운전 영상은 서울 6개, 도쿄 5개인데 부산 3개, 오사카 1개, 타이베이 0개였다.

이에 따라 b6를 만들었다. 세 도시에 서로 다른 채널의 운전 영상을 4\textasciitilde 5개씩 넣어 서울·도쿄와 비슷하게 맞췄다. 영상을 너무 많이 넣으면 이번에는 서울과 도쿄가 뒤집힐 수 있어, 더한 뒤에도 영상 수는 서울(42)과 도쿄(46)보다 적게 했다(타이베이 25, 부산 33, 오사카 38). 오사카에는 도쿄에 비슷한 거리가 없는 도톤보리의 낮 영상을 넣었다. 밤의 도톤보리는 오답이 거의 없었기 때문이다.

\paragraph{고른 영상은 정말 부족한 장면이었다.} 새로 넣은 영상의 프레임을, 그 영상을 본 적 없는 b5는 오사카 31\%, 부산 32\%, 타이베이 52\%만 맞혔고, 틀린 프레임은 오사카는 도쿄로, 부산은 서울로 갔다. 오답을 보고 고른 영상이 학습 데이터에 없던 장면을 정확히 채운 것이다.

\paragraph{결과.} 새 검증은 0.486에서 0.523으로 올랐다(표~\ref{tab:newval}). 새 검증 사진을 걷기 영상의 프레임과 운전 영상의 프레임으로 나눠 보면, 오른 것은 주로 운전 영상의 프레임이었다(표~\ref{tab:walkdrive}). 처음 보는 운전 영상의 프레임 정답률이 오사카 .118→.314, 타이베이 .367→.617, 부산 .333→.467로 크게 올랐다.

\begin{table}[h]\centering\footnotesize
\begin{tabular}{lrrrrrrr}
\toprule
모델 & 공식 검증 & 새 검증 & 서울 & 부산 & 도쿄 & 오사카 & 타이베이 \\
\midrule
b5 & 0.580 & 0.486 & .464 & .507 & .614 & .243 & .600 \\
b6 & 0.579 & 0.523 & .436 & .571 & .479 & .429 & .700 \\
b7 & 0.582 & 0.513 & .407 & .621 & .436 & .443 & .657 \\
\bottomrule
\end{tabular}
\caption{b5, b6, b7의 새 검증 도시별 정답률 (\texttt{outputs/new\_val\_b5\_b6\_b7.txt}).}
\label{tab:newval}
\end{table}

\begin{table}[h]\centering\footnotesize
\begin{tabular}{lll}
\toprule
도시 & 걷기 영상의 프레임 (b5 → b6 → b7) & 운전 영상의 프레임 (b5 → b6 → b7) \\
\midrule
오사카 & .315 → .494 → .539 & .118 → .314 → .275 \\
타이베이 & .775 → .762 → .762 & .367 → .617 → .517 \\
부산 & .637 → .650 → .688 & .333 → .467 → .533 \\
\bottomrule
\end{tabular}
\caption{새 검증에서 걷기 영상과 운전 영상의 프레임을 나눠 잰 정답률 (LOG의 b6, b7 항목).}
\label{tab:walkdrive}
\end{table}

새 검증 영상은 학습 영상과 채널이 겹치지 않으므로, 영상 하나를 외운 것이 아니라 그 도시의 도로 모습을 어느 정도 배운 것으로 예상한다.

새 검증의 오사카 정답률은 .243에서 .429로 다섯 도시 중 가장 크게 올랐다. 평범한 동네 걷기 영상은 1개만 넣어서, 오답의 대부분인 평범한 거리 사진은 덜 고쳐질 것으로 예상했었다. 오사카가 오른 이유는 두 가지로 나눠 볼 수 있다.
\begin{itemize}
\item \textbf{도쿄로 쏠리던 답이 줄었다.} 새 검증의 모든 답 중 도쿄로 답한 비율은 .287에서 .216으로 줄고, 오사카로 답한 비율은 .107에서 .213으로 늘었다. 도쿄 정답률이 .614에서 .479로 내려간 것을 더해도, 도쿄와 오사카 정답률의 합은 .857에서 .908로 올랐다.
\item \textbf{도톤보리 사진을 맞히게 됐다.} b5의 오답에 여러 장 있던 도톤보리 강가 사진들이 b6의 오답에서는 거의 사라졌다. 오사카 거리를 더 배운 것도 점수 상승에 기여한 것으로 예상된다.
\end{itemize}
다만 모델이 도쿄 사진을 오사카로 답하는 비율도 .221에서 .429로 늘었다. 일본 사진이 이제 도쿄와 오사카에 거의 반반씩 가는 셈이라, 평범한 거리에서 두 도시를 가르는 힘은 아직 약하다. 서울 사진의 40\%를 부산으로 답하는 것도 같은 문제다. 공식 검증은 0.580에서 0.579로 거의 그대로였다.

\subsection*{7. 학습 방법만 바꾸기 (b7)}

마지막으로 데이터는 그대로 두고 학습 방법만 바꿨다. b6의 오답은 대부분 나도 간판의 작은 글자나 도로와 건물의 디테일을 보지 않으면 구분하지 못할 사진이었다. 거리 사진을 2배 자주 보여 주면 이런 작은 단서까지 배울 것을 기대했고, 배운다면 눈에 띄게 오르고, Track A처럼 외우기만 늘면 조금 떨어질 것으로 예상했다. 결과는 공식 검증 +0.3\%p, 새 검증 −1.0\%p로 후자에 가까웠다(표~\ref{tab:newval}). 다만 두 차이 모두 같은 설정을 다시 학습할 때의 흔들림과 구별하기 어려운 크기다. 도시별로는 부산이 5.0\%p 올랐지만 서울, 도쿄, 타이베이는 2.9\textasciitilde 4.3\%p 내려가 방향이 엇갈렸고, 타이베이 운전 영상의 프레임 정답률은 .617에서 .517로 내려갔다(표~\ref{tab:walkdrive}). b6에서도 학습의 28\textasciitilde 50\%가 거리 사진이었으므로 대강의 단서는 b6에서 이미 배웠고, 더 디테일한 단서는 배우지 못한 것으로 예상한다. b6의 오답이 대부분 그럴듯했다는 점이 이를 뒷받침한다. Track A에서도 거리 사진 가중치는 효과가 없었고, Track B에서도 학습 방법보다 부족한 장면의 영상을 더하는 쪽이 효과가 컸다. b7은 릴리스하지 않았다.

\subsection*{검증 점수와 리더보드가 어긋나는 이유}

\begin{center}\footnotesize
\begin{tabular}{lrrrrrrr}
\toprule
모델 (릴리스) & 공식 검증 & 새 검증 & 리더보드 & 서울 & 부산 & 도쿄 & 오사카 / 타이베이 \\
\midrule
b2 (format) & 0.585 & 0.474 & 0.432 & .61 & .33 & .74 & .30 / .17 \\
b5 (seoul-tokyo) & 0.580 & 0.486 & 0.468 & .64 & .36 & .75 & .37 / .22 \\
b6 (drive-dotonbori) & 0.579 & 0.523 & 0.478 & .65 & .40 & .70 & .40 / .23 \\
\bottomrule
\end{tabular}
\end{center}

b6의 리더보드는 b5의 0.468에서 0.478로 1.0\%p 올랐다. 도쿄가 .75에서 .70으로 조금 내려간 것을 빼면, 나머지 네 도시가 모두 올랐다(서울 .64→.65, 부산 .36→.40, 오사카 .37→.40, 타이베이 .22→.23). 특히 서울과 부산이 동시에 올랐다. 한 도시가 오르면 같은 나라의 다른 도시가 내려가던 이전과 달리, 두 도시를 가르는 단서를 더 잘 학습한 것으로 보인다. 미리 정한 대로 b6가 b5보다 높으므로 b6를 최종으로 두었고, b6는 Track B 리더보드에서 1위였다.

타이베이는 미스테리로 남았다. 새 검증에서는 늘 높았고, 학습 영상도 밤, 비, 운전 등 다양하게 더했는데 리더보드는 .22에서 .23으로 거의 그대로다. 타이베이의 새 검증, 공식 검증과 리더보드 사진이 다른 것은 확실하다. 하지만 학습 사진을 다양하게 더해 학습이 개선되었다면, 새 검증 사진과 달라도 리더보드 점수가 더 높아야 한다고 생각한다. 리더보드 타이베이 점수는 모든 학생이 낮고(.11\textasciitilde .29), 224 사진을 쓰는 Track A 모델(.21\textasciitilde .29)도 448 모델보다 낮지 않다. 리더보드의 타이베이 사진은 우리가 더한 장면과 다른 종류인 것으로 예상되지만, 리더보드 사진은 볼 수 없어 확인할 수 없다. 검증 점수와 리더보드가 어긋나는 이유를 정리하면 다음과 같다.

\begin{itemize}
\item \textbf{사진의 종류가 다르다.} 공식 검증의 약 70\%는 위키미디어 사진이지만 리더보드는 모두 거리 영상이다. 그래서 공식 검증은 거리 영상을 더한 효과를 거의 보여 주지 못했다(b5→b6에서 0.580→0.579).
\item \textbf{검증 영상이 적다.} 새 검증도 도시마다 영상이 7\textasciitilde 8개뿐이다. 그 몇 개가 어떤 장면이냐에 따라 도시 점수가 크게 달라진다. 그래서 검증은 어느 도시가 어려운지는 맞히지만(r = 0.81), 모델을 바꿨을 때 도시별 변화의 방향은 맞히지 못했다(r = −0.11). b5에서 서울과 도쿄가 그 예다.
\item \textbf{도시마다 리더보드 사진의 구성이 다른 것으로 예상된다.} 타이베이는 모든 검증에서 리더보드보다 훨씬 높다. 우리 학습과 검증의 타이베이 영상과 리더보드의 타이베이 영상이 다른 장면(밤, 운전 등)으로 이루어져 있는 것으로 예상하지만, 리더보드 사진은 볼 수 없어 직접 확인할 수는 없다.
\end{itemize}

\subsection*{제출 파일 확인 (채점하는 쪽이 보는 것)}
\outputfile{outputs/check_submission.txt}

% ---------------------------------------------------------------------------
\section*{무엇이 차이를 만들었나}

\paragraph{가장 도움이 된 변경: 다섯 도시에 같은 조건으로, 부족한 장면의 영상을 더한 것.}
리더보드에서 흔들림보다 큰 차이를 낸 유일한 변경이다. b5는 b2보다 리더보드가 3.6\%p 올랐고 다섯 도시가 모두 같은 방향으로 움직였다. b6는 여기서 1.0\%p 더 올라 0.478로 Track B 1위가 되었다. Track A의 변경은 모두 1\textasciitilde 3\%p였고, 같은 레시피를 seed만 바꿔도 4.5\%p가 움직였다. 이 효과는 영상을 넣기만 해서 생긴 것이 아니었다. 세 도시에만 넣은 b3, b4는 지름길을 배웠고, 화질과 촬영 스타일을 다섯 도시에 같게 맞춘 뒤에야(b5) 효과가 났다. b6에서는 오답을 영상별로 나눠 부족한 장면(운전, 도톤보리)을 찾았고, 그 영상을 기존 모델이 31\textasciitilde 52\%밖에 맞히지 못한다는 것도 확인했다. 즉 효과를 낸 것은 ``데이터를 많이 넣는 것''이 아니라 ``모델이 모르는 장면을 찾아, 도시 사이의 조건을 맞춰 넣는 것''이었다. 이는 Track A에서 내린 결론, 즉 낮은 점수의 원인이 외우기보다 학습 영상의 종류와 수라는 예상과 맞는다.

공식 검증에서 가장 크게 오른 것은 b1→b2(모델, 448, 도시 균등, +9.2\%p)였다. 글자를 단서로 쓰게 된 것도 이때다. 다만 세 가지를 한 번에 바꿔 각각의 기여는 가를 수 없다.

\paragraph{가장 도움이 안 된 변경.}
사진 정제(Track A 1\textasciitilde 3차)와 거리 사진 가중치(Track A clean\_street, Track B b7)는 점수를 흔들림 이상으로 움직이지 못했다. 층별 학습률(−1.6\%p), crop60(−1.6\%p), 흐림(−1.7\%p)은 공식 검증을 내렸다. 외우기를 줄이는 방법(mixup, CutMix)은 외우기는 실제로 줄였지만 처음 보는 영상의 점수는 올리지 못했다. 있는 사진을 빼거나 가리거나 더 자주 보여 주는 것은, 모델이 본 적 없는 장면을 알려 주지 못했다.

\paragraph{일주일이 더 있다면.}
\begin{itemize}
\item \textbf{글자를 확실하게 구분하는 모델.} b6의 오답을 보면 거리 모습만으로는 나도 구분하기 힘든 사진이 꽤 많다. 이런 사진에서 나라를 확실히 가를 수 있는 단서는 언어다. 지금 모델은 글자가 크게 쓰여 있으면 구분하는 것 같지만, 글씨체가 독특하거나 두껍거나 글자가 작으면 구분하지 못한다. 다른 나라로 답한 오답의 대부분이 이런 경우였다. 예를 들어 두꺼운 한글 간판을 보고 타이베이로 답했고, 새 검증에서는 번체 한자 간판이 화면 가득한 타이베이 가게 앞 사진을 오사카로 답했다. 글씨를 못 읽은 것이다. 이런 글자는 거리 사진을 더 많이 보여 주는 것만으로는 배우기 어렵다. 글씨체가 바뀌면 글자의 대략적인 생김새도 바뀌며, 거리 사진의 다른 단서들을 주요하게 배울 가능성도 높기 때문이다. 그래서 여러 글씨체와 굵기, 크기로 그린 한글, 가나, 번체 한자 이미지로 먼저 글자 종류를 구분하도록 학습하고, 그 모델 위에서 도시 분류를 학습하고 싶다. 모델 크기 제한 안에서 할 수 있다. 다만 번체 한자와 일본 한자는 같은 글자가 많아 한글보다 어려울 것으로 예상된다.

b2가 글자를 단서로 쓰는지 확인한 실험(Track B 1절)에서는, 언어로 도시를 가를 수 있는 사진이 적게 나왔다. 이 실험에서 OCR(사진 속 글자를 찾는 프로그램)이 그 나라 글자(한글, 가나, 한자)를 3군데 이상에서 찾은 사진은 공식 검증 사진 11,758장 중 228장(약 2\%)뿐이었다. 그래서 같은 실험에서 글자를 회색으로 가렸을 때 b2의 전체 정답률도 1\%p만 떨어졌다(59.9\% → 58.9\%). 그러나 이것은 OCR 모델이 애초에 글자를 잘 찾고 분류하지 못했기 때문이기도 하다. OCR이 분류한 결과를 보면 틀린 것도 많았고, 나라와 맞지 않는 글자가 나온 사진 319장(한국 밖 사진의 한글, 일본 밖 사진의 가나)은 아예 빼야 했다. 그래서 언어로 분류할 수 있는 도시 사진은 이 실험의 수치보다 더 많을 것으로 생각한다. 글자 구분이 확실해지면 효과도 그만큼 클 것으로 예상한다.
\item \textbf{리더보드를 닮은 검증과 타이베이.} 타이베이는 새 검증에서 0.700인데 리더보드는 .23이다. b6의 새 검증 타이베이 오답 42장 중 25장이 밤 시내 운전 영상 하나(14장)와 현대 도심 걷기 영상 하나(11장)에서 나왔다. 학습에 넣은 타이베이 운전 영상은 흐린 낮의 고속도로가 대부분이라, 밤 시내 도로와 고층 빌딩이 있는 현대 도심은 거의 배우지 못한 것으로 예상된다. 그래서 먼저 검증을 바꾸고 싶다. 도시마다 영상 수를 늘리고 영상당 장수를 줄여 영상 하나가 점수를 흔들지 않게 하고, 낮과 밤, 옛 시가지와 현대 도심, 걷기와 운전이 골고루 들어가게 한다. 다섯 도시를 같은 방식으로 바꾸고, 리더보드처럼 프레임은 거르지 않는다. 바꾼 검증으로 리더보드 점수가 있는 b2, b5, b6를 다시 재어 리더보드와 잘 맞는지 확인한 뒤, 그 검증으로 타이베이에 부족한 장면(밤 시내 운전, 현대 도심 걷기)의 영상을 골라 넣고 싶다.
\end{itemize}

% ---------------------------------------------------------------------------
\section*{All runs}
\outputfile{outputs/experiments.csv}

\section*{Lab notebook}
\outputfile{outputs/LOG.md}

\end{document}
